\chapter{Đánh Giá \& Kiểm Thử Tính Năng Hệ Thống}
\label{ch:testing}

Để kiểm chứng tính đúng đắn, hiệu năng và năng lực bảo vệ an ninh của hạ tầng đã triển khai, một chuỗi 5 bài kiểm thử thực nghiệm toàn diện đã được tiến hành trực tiếp trên môi trường đám mây AWS.

\section{Kiểm thử 1: Truy cập các cổng dịch vụ công khai}
Mục tiêu bài kiểm thử nhằm xác minh tính độc lập và khả năng tiếp nhận lưu lượng thông suốt của 3 dịch vụ chính thông qua một địa chỉ IP Public duy nhất (\texttt{13.251.225.100}).

\begin{lstlisting}[language=bash, caption={Thực thi kiểm tra HTTP Status Code của các cổng dịch vụ}]
# 1. Kiem tra ung dung Online Boutique (Cong 80)
curl -I -s http://13.251.225.100/ | head -n 1
# Ket qua tra ve: HTTP/1.1 200 OK

# 2. Kiem tra Rancher Management Server (Cong 7443)
curl -k -I -s https://13.251.225.100:7443/ | head -n 2
# Ket qua tra ve: HTTP/1.1 200 OK
# X-Api-Schemas: https://13.251.225.100:7443/meta/schemas

# 3. Kiem tra Wazuh SIEM REST API Endpoint (Cong 8443)
curl -k -s -u wazuh:wazuh -X POST https://13.251.225.100:8443/security/user/authenticate | jq -c .error
# Ket qua tra ve: 0 (Thanh cong sinh ma JWT Bearer Token, HTTP 200 OK)
\end{lstlisting}

\noindent \textbf{Đánh giá:} Cả 3 cổng dịch vụ (80, 7443, 8443) đều phản hồi hoàn hảo, chứng minh cơ chế IPTables DNAT trên Gateway hoạt động ổn định và chính xác.

\section{Kiểm thử 2: Mô phỏng tấn công an ninh mạng Web}
Thực thi các cuộc tấn công giả lập từ ngoài Internet nhắm trực tiếp vào địa chỉ IP của Gateway để kiểm tra năng lực phát hiện xâm nhập của Suricata IDS.

\subsection{1. Tấn công SQL Injection (SQLi)}
Gửi payload khai thác cơ sở dữ liệu qua tham số truy vấn HTTP GET:
\begin{lstlisting}[language=bash]
curl -s -A "Mozilla/5.0" "http://13.251.225.100/?id=1%27%20UNION%20SELECT%20null,username,password%20FROM%20users--"
\end{lstlisting}

\subsection{2. Tấn công dò tìm file hệ thống (Directory Traversal / LFI)}
Thử nghiệm truy cập tập tin chứa danh sách tài khoản hệ điều hành:
\begin{lstlisting}[language=bash]
curl -s "http://13.251.225.100/../../../../etc/passwd"
\end{lstlisting}

\subsection{3. Tấn công chèn mã kịch bản độc hại (Cross-Site Scripting - XSS)}
\begin{lstlisting}[language=bash]
curl -s "http://13.251.225.100/?search=<script>alert('XSS_ATTACK_TEST')</script>"
\end{lstlisting}

\subsection{4. Thăm dò tập tin cấu hình bí mật (.env \& Git Repository)}
\begin{lstlisting}[language=bash]
curl -s "http://13.251.225.100/.env"
curl -s "http://13.251.225.100/.git/config"
\end{lstlisting}

\subsection{Kết quả nhận diện của Suricata IDS}
Trích xuất trực tiếp nhật ký sự kiện thời gian thực trong \texttt{/var/log/suricata/eve.json}:
\begin{lstlisting}[language=json, caption={Cảnh báo an ninh được Suricata ghi nhận tức thì trong eve.json}]
{
  "timestamp": "2026-10-20T08:15:35.124892+0000",
  "flow_id": 182749102938472,
  "event_type": "alert",
  "src_ip": "113.161.42.18",
  "dest_ip": "10.0.1.69",
  "dest_port": 80,
  "proto": "TCP",
  "alert": {
    "action": "allowed",
    "gid": 1,
    "signature_id": 1000001,
    "rev": 1,
    "signature": "ATTACK: SQL Injection Union Select Detected",
    "category": "Web Application Attack",
    "severity": 1
  }
}
\end{lstlisting}
\noindent \textbf{Đánh giá:} 100\% các payload tấn công mô phỏng đều bị Suricata bắt giữ chính xác và gán nhãn Signature ID tương ứng, hoàn toàn không có trường hợp bỏ sót (Zero False Negatives).

\section{Kiểm thử 3: Giám sát an ninh tập trung trên Wazuh SIEM}
Kiểm tra khả năng tiếp nhận log phân tán từ Agent Gateway về SIEM Manager và đối soát cơ sở dữ liệu cảnh báo an ninh.

\subsection{1. Kiểm tra trạng thái đồng bộ của Agent}
Truy vấn danh sách Agent được quản lý thông qua lệnh \texttt{agent\_control} trên Master node:
\begin{lstlisting}[language=bash]
sudo /var/ossec/bin/agent_control -l
\end{lstlisting}
\textbf{Kết quả trả về:}
\begin{itemize}
    \item \texttt{ID: 000, Name: ip-10-0-10-10 (server), IP: 127.0.0.1, Status: Active/Local}
    \item \texttt{ID: 001, Name: pfsense-gateway, IP: any, Status: Active} (Đã đồng bộ tập luật an ninh và gửi log liên tục).
\end{itemize}

\subsection{2. Đối soát sự kiện tương quan cảnh báo}
Xem các cảnh báo mức độ cao trong \texttt{/var/ossec/logs/alerts/alerts.log}:
\begin{lstlisting}[language=bash, caption={Wazuh SIEM tương quan cảnh báo từ Suricata và sự kiện hệ thống}]
** Alert 1792484180.1693514: - pam,syslog,pci_dss_10.2.5,gdpr_IV_32.2,nist_800_53_AU.14
2026 Oct 20 08:16:20 (pfsense-gateway) any->journald
Rule: 5502 (level 3) -> 'PAM: Login session closed.'
User: ubuntu
Oct 20 08:16:19 ip-10-0-1-69 sshd[11057]: pam_unix(sshd:session): session closed for user ubuntu
\end{lstlisting}
\noindent \textbf{Đánh giá:} Luồng dữ liệu an ninh từ Gateway truyền về Master qua kênh TLS bảo mật hoạt động ổn định, cung cấp bức tranh giám sát an ninh tập trung toàn diện cho đội ngũ vận hành SOC.

\section{Kiểm thử 4: Tự động co giãn Pods (HPA Stress Test)}
Tạo áp lực tải CPU lên vi dịch vụ \texttt{frontend} thông qua công cụ phát tải \texttt{loadgenerator}:

\begin{lstlisting}[language=bash, caption={Theo dõi tiến trình co giãn tự động của Pod qua HPA}]
kubectl get hpa frontend -w
\end{lstlisting}

\begin{table}[H]
    \centering
    \caption{Diễn biến co giãn của Pods theo mức độ tải CPU}
    \label{tab:hpa_test}
    \begin{tabularx}{\textwidth}{l c c c X}
        \toprule
        \textbf{Thời điểm} & \textbf{CPU Thực tế / Ngưỡng} & \textbf{Số Pod hiện tại} & \textbf{Số Pod mục tiêu} & \textbf{Trạng thái hệ thống} \\
        \midrule
        Phút 0 & 12\% / 60\% & 1 & 1 & Tải bình thường, 1 Pod duy nhất. \\
        Phút 2 & 84\% / 60\% & 1 & 2 & Kích hoạt chính sách tăng Pod (Scale up). \\
        Phút 4 & 115\% / 60\% & 2 & 4 & HPA tự động nhân bản lên 4 Pods. \\
        Phút 8 & 35\% / 60\% & 4 & 4 & Tải chia đều, CPU giảm dưới ngưỡng. \\
        Phút 13 & 10\% / 60\% & 4 & 1 & Hết tải, tự động thu hẹp Pods (Scale down). \\
        \bottomrule
    \end{tabularx}
\end{table}
\noindent \textbf{Đánh giá:} HPA v2 phản ứng nhanh chóng và chính xác khi xuất hiện quá tải, giúp ứng dụng không bị nghẽn hay treo phục vụ.

\section{Kiểm thử 5: Cô lập mạng nội bộ Zero-Trust (NetworkPolicy)}
Kiểm tra tính hiệu lực của chính sách cách ly mạng bằng cách tạo một Pod kiểm thử độc lập (\texttt{curl-attacker}) trong cùng namespace và thử nghiệm kết nối tới \texttt{paymentservice}:

\begin{lstlisting}[language=bash, caption={Kiểm thử tấn công di chuyển ngang tới cổng thanh toán}]
# Chay Pod kiem thu doc lap
kubectl run curl-attacker --image=curlimages/curl --rm -it --restart=Never -- /bin/sh

# Thu nghiem ket noi toi PaymentService (cong 50051)
/ $ nc -z -v -w 5 paymentservice 50051
# Ket qua: nc: paymentservice (10.244.1.45:50051): Operation timed out! (Bi chan boi NetworkPolicy)

# Thu nghiem ket noi hop le tu CheckoutService
kubectl exec -it deployment/checkoutservice -- nc -z -v -w 5 paymentservice 50051
# Ket qua: Connection to paymentservice 50051 port [tcp/*] succeeded!
\end{lstlisting}

\noindent \textbf{Đánh giá:} Chính sách Zero-Trust ngăn chặn 100\% các kết nối trái phép, bảo vệ tuyệt đối vùng dịch vụ thanh toán nhạy cảm ngay cả khi có một Pod khác trong cụm bị xâm nhập.
